% =============================================================
% บทที่ 2 - ทฤษฎีและงานวิจัยที่เกี่ยวข้อง
% โครงงานประเภท B (งานวิจัย) — Algorithm, Statistical Methods,
% Dataset, Evaluation Metrics + ตารางเปรียบเทียบ + Research Gap
% =============================================================

% -----------------------------------------------------------
\section{ทฤษฎีที่เกี่ยวข้อง}
% -----------------------------------------------------------

\subsection{การสร้างคำตอบโดยอาศัยการค้นคืน}

การสร้างคำตอบโดยอาศัยการค้นคืน หรือ Retrieval-Augmented Generation (RAG) เป็นวิธีที่ให้แบบจำลองภาษาตอบคำถามโดยอ้างอิงเอกสารที่ค้นมาได้ แทนที่จะตอบจากความรู้ที่ฝังอยู่ในตัวแบบจำลองเพียงอย่างเดียว \cite{lewis2020rag} วิธีนี้แก้ปัญหาสองข้อพร้อมกัน คือแบบจำลองไม่รู้ข้อมูลที่เกิดขึ้นหลังการฝึก และแบบจำลองมักแต่งข้อมูลขึ้นเองเมื่อไม่รู้คำตอบ

โครงสร้างพื้นฐานของ RAG ทำงานเป็นท่อเส้นตรง คือรับคำถาม ค้นหนึ่งครั้ง แล้วเขียนคำตอบจากสิ่งที่ค้นได้ จุดอ่อนที่ตามมาคือถ้าการค้นครั้งเดียวนั้นหาของไม่เจอ ระบบไม่มีโอกาสแก้ตัว และยังตอบด้วยของที่ไม่ตรงอยู่ดี \cite{gao2023survey}

\subsection{การแบ่งประเภทของสถาปัตยกรรม RAG}

งานทบทวนวรรณกรรมแบ่งสถาปัตยกรรม RAG ออกเป็นหลายระดับ \cite{gao2024modular} โดยนับการค้นซ้ำแบบวนรอบ การค้นแบบเรียกตัวเอง และการค้นที่ปรับตามความยากของคำถาม ว่าอยู่ในกลุ่ม RAG ทั้งหมด เพราะลำดับการทำงานยังถูกกำหนดไว้ล่วงหน้าด้วยโปรแกรม

ส่วน Agentic RAG นิยามด้วยเกณฑ์ที่ต่างออกไป คือต้องมีชั้นที่แบบจำลองภาษาเป็นผู้ตัดสินใจเองขณะทำงาน ไม่ใช่เพียงทำตามลำดับที่เขียนไว้ \cite{singh2025agentic} รูปแบบที่ง่ายที่สุดเรียกว่า single-agent router คือมีตัวแทนหนึ่งตัวทำหน้าที่เลือกว่าจะทำอะไรต่อ เส้นแบ่งนี้สำคัญกับโครงงานนี้ เพราะเป็นเกณฑ์ที่ใช้จัดประเภทให้กับสถาปัตยกรรมทั้งสามแบบที่นำมาเปรียบเทียบ

\textbf{ข้อสรุปที่ได้จากวรรณกรรม} คือการมีลูปค้นซ้ำไม่ได้ทำให้ระบบเป็น Agentic RAG สิ่งที่ทำให้เป็นคือการมีชั้นที่แบบจำลองภาษาตัดสินใจเองขณะทำงาน

\subsection{การฝังเวกเตอร์และการค้นคืนเชิงความหมาย}

การค้นคืนเชิงความหมายแปลงทั้งคำค้นและเอกสารเป็นเวกเตอร์ในปริภูมิเดียวกัน แล้ววัดความใกล้เคียงด้วยค่าโคไซน์ ทำให้ค้นเจอแม้ใช้คำคนละคำกับที่เอกสารใช้ ซึ่งการค้นแบบจับคู่คำตรงตัวทำไม่ได้

โครงงานนี้ใช้แบบจำลองฝังเวกเตอร์ที่รองรับหลายภาษาในแบบจำลองเดียว \cite{chen2024bgem3} จึงวางคำค้นภาษาไทยและชื่อสินค้าภาษาอังกฤษไว้ในปริภูมิเดียวกันได้ ส่วนการค้นเวกเตอร์ที่ใกล้ที่สุดในคลังขนาดใหญ่ใช้โครงสร้างดัชนีแบบกราฟลำดับชั้น \cite{malkov2020hnsw} ซึ่งแลกความแม่นยำสมบูรณ์กับความเร็วที่เพิ่มขึ้นหลายเท่า

\subsection{การจัดอันดับด้วยการรวมอันดับจากหลายสัญญาณ}

เมื่อมีสัญญาณหลายตัวที่บอกว่าสินค้าชิ้นใดควรอยู่อันดับต้น เช่น ความใกล้เชิงความหมาย คะแนนดาว และจำนวนรีวิว การรวมสัญญาณด้วยการบวกค่าดิบทำได้ยาก เพราะแต่ละสัญญาณมีหน่วยและช่วงค่าต่างกัน

วิธีรวมอันดับแบบ Reciprocal Rank Fusion \cite{cormack2009rrf} แก้ปัญหานี้โดยไม่ใช้ค่าดิบเลย แต่ใช้เฉพาะ\textbf{อันดับ}ที่แต่ละสัญญาณจัดให้ แล้วรวมคะแนนตามสมการที่ \ref{eq:rrf}

\begin{equation}
\label{eq:rrf}
\mathrm{RRF}(d) = \sum_{i \in S} \frac{w_i}{k + r_i(d)}
\end{equation}

เมื่อ $d$ คือสินค้าที่พิจารณา $S$ คือเซตของสัญญาณทั้งหมด $w_i$ คือน้ำหนักของสัญญาณที่ $i$ $r_i(d)$ คืออันดับของสินค้านั้นตามสัญญาณที่ $i$ และ $k$ คือค่าคงที่ที่ลดอิทธิพลของอันดับต้น ๆ ไม่ให้มากเกินไป

คุณสมบัติที่ต้องเข้าใจคือ RRF ใช้เฉพาะอันดับ จึง\textbf{ไม่รับรู้ว่าสัญญาณหนึ่งเหนือกว่าอีกสัญญาณมากเพียงใด} คุณสมบัตินี้มีผลโดยตรงต่อข้อค้นพบของโครงงานนี้

\subsection{ตัวชี้วัดคุณภาพการค้นคืน}

ตัวชี้วัดหลักที่ใช้คือ Normalized Discounted Cumulative Gain (nDCG) \cite{jarvelin2002ndcg} ซึ่งให้คะแนนตามความตรงของสิ่งที่ค้นได้ โดยถ่วงน้ำหนักให้อันดับต้นมีค่ามากกว่าอันดับท้าย ตามสมการที่ \ref{eq:dcg} และ \ref{eq:ndcg}

\begin{equation}
\label{eq:dcg}
\mathrm{DCG@}k = \sum_{i=1}^{k} \frac{g_i}{\log_2(i+1)}
\end{equation}

\begin{equation}
\label{eq:ndcg}
\mathrm{nDCG@}k = \frac{\mathrm{DCG@}k}{\mathrm{IDCG@}k}
\end{equation}

เมื่อ $g_i$ คือค่าความตรงของสินค้าที่อันดับ $i$ และ $\mathrm{IDCG@}k$ คือค่าของการจัดอันดับที่ดีที่สุดเท่าที่เป็นไปได้ ซึ่งใช้เป็นตัวหารเพื่อให้คะแนนอยู่ในช่วง 0 ถึง 1

ตัวชี้วัดนี้มีข้อจำกัดสำคัญคือ\textbf{วัดเฉพาะคุณภาพของสิ่งที่ตอบ แต่วัดการปฏิเสธที่ถูกต้องไม่ได้} ระบบที่ตอบว่าไม่พบสินค้าในกรณีที่คลังไม่มีของจริง จะได้ nDCG เท่ากับศูนย์เช่นเดียวกับระบบที่ตอบของผิด ทั้งที่พฤติกรรมทั้งสองต่างกันโดยสิ้นเชิง โครงงานนี้จึงต้องเพิ่มตัวชี้วัดการตัดสินใจแนะนำหรือปฏิเสธแยกต่างหาก

\subsection{ความซื่อตรงของคำตอบ}

นอกจากคำถามว่าระบบหาของตรงไหม ยังมีคำถามว่าระบบเขียนเกินข้อมูลที่มีอยู่ในมือหรือไม่ ปัญหานี้เรียกรวมว่าการแต่งข้อมูล ซึ่งเป็นปัญหาที่พบทั่วไปในแบบจำลองภาษา \cite{ji2023hallucination}

งานวิจัยแบ่งความผิดพลาดเป็นสองชนิด \cite{maynez2020faithfulness} คือ intrinsic ซึ่งเป็นการบิดเบือนข้อมูลที่ให้ไป และ extrinsic ซึ่งเป็นการเติมสิ่งที่ไม่มีอยู่ในข้อมูลเลย หลักที่ใช้ตัดสินคือทุกข้อความในคำตอบต้องมีข้อมูลรองรับ \cite{rashkin2023attribution} โครงงานนี้ใช้การแบ่งชนิดนี้ในการติดป้ายความซื่อตรงของคำตอบ

\subsection{การสร้างชุดทดสอบและการวัดความน่าเชื่อถือของเฉลย}

การประเมินระบบค้นคืนใช้กรอบ Cranfield ซึ่งประกอบด้วยคลังเอกสาร ชุดคำค้น และเฉลยความตรง \cite{voorhees2005trec} การเลือกคู่ที่จะติดเฉลยนิยมใช้วิธี pooling คือรวมสิ่งที่ระบบต่าง ๆ หยิบขึ้นมาเป็นกองเดียวแล้วติดเฉลยครั้งเดียว

ข้อจำกัดที่ต้องระวังคือ pooling bias \cite{zobel1998reliable} กล่าวคือกองที่ติดเฉลยย่อมขึ้นกับระบบที่สร้างมันขึ้นมา ระบบใหม่ที่หยิบของนอกกองจะถูกประเมินต่ำกว่าความเป็นจริง ข้อจำกัดนี้เกิดขึ้นจริงในโครงงานนี้เมื่อเปลี่ยนชั้นจัดอันดับ และเป็นเหตุให้ต้องติดเฉลยเพิ่ม

ส่วนความน่าเชื่อถือของเฉลยวัดด้วยค่าความสอดคล้องระหว่างผู้ติดป้าย โดยใช้ weighted kappa \cite{cohen1968kappa} ซึ่งลงโทษความต่างมากขึ้นตามระยะห่างของระดับป้าย อ่านผลด้วยช่วงมาตรฐาน \cite{landis1977agreement} และรายงานตามแนวปฏิบัติของงานด้านภาษา \cite{artstein2008agreement} ส่วนกระบวนการติดป้ายที่ดีควรเขียนคู่มือก่อน ติดกองนำร่อง แล้วจึงติดทั้งกอง \cite{pustejovsky2012annotation}

\subsection{วิธีทางสถิติที่ใช้ทดสอบนัยสำคัญ}

เมื่อเปรียบเทียบระบบสองระบบบนชุดคำค้นเดียวกัน ข้อมูลที่ได้เป็นข้อมูลจับคู่ คือคำค้นเดียวกันถูกวัดด้วยทั้งสองระบบ การทดสอบที่เหมาะสมกับข้อมูลลักษณะนี้และไม่ต้องสมมติว่าข้อมูลแจกแจงปกติ คือการทดสอบแบบสลับเครื่องหมาย \cite{smucker2007comparison}

หลักการคือนำผลต่างของแต่ละคำค้นมาสลับเครื่องหมายทุกรูปแบบที่เป็นไปได้ แล้วนับว่ามีกี่รูปแบบที่ให้ผลรวมมากกว่าหรือเท่ากับค่าที่สังเกตได้จริง สัดส่วนนั้นคือค่า $p$ ข้อควรระวังคือเมื่อจำนวนคำค้นที่ให้ผลต่างกันมีน้อย ค่า $p$ จะมีขีดจำกัดล่างที่ไม่สามารถต่ำกว่านั้นได้ เช่น เมื่อมีคำค้นที่ต่างกันเพียง 5 ข้อ ค่า $p$ ต่ำสุดที่เป็นไปได้คือ $2/2^5 = 0.0625$

ส่วนขนาดของชุดทดสอบ งานวิจัยด้านการประเมินแนะนำว่าจำนวนคำค้นมีผลต่อความเสถียรของผลมากกว่าความละเอียดของเฉลย \cite{buckley2000stability, voorhees2000variations}

% -----------------------------------------------------------
\section{งานวิจัยที่เกี่ยวข้อง}
% -----------------------------------------------------------

\subsection{งานที่วางรากฐานของ RAG}

Lewis และคณะ \cite{lewis2020rag} เสนอสถาปัตยกรรม RAG เป็นครั้งแรก โดยรวมตัวค้นคืนเข้ากับแบบจำลองสร้างข้อความ และแสดงว่าวิธีนี้ให้คำตอบที่ถูกต้องกว่าการใช้แบบจำลองเพียงอย่างเดียวในงานที่ต้องใช้ความรู้ \textbf{ส่วนที่คล้ายกับโครงงานนี้} คือใช้โครงค้นแล้วตอบเหมือนกัน \textbf{ส่วนที่ต่าง} คืองานนั้นค้นครั้งเดียวและไม่มีชั้นตัดสินใจ ซึ่งในโครงงานนี้ใช้เป็นเส้นฐานเปรียบเทียบ \textbf{ข้อจำกัด} คือเมื่อการค้นครั้งเดียวล้มเหลว ระบบไม่มีกลไกแก้ตัว

\subsection{งานที่เพิ่มความสามารถตัดสินใจให้ระบบ RAG}

Jeong และคณะ \cite{jeong2024adaptive} เสนอวิธีให้ระบบเลือกกลยุทธ์การค้นตามความยากของคำถาม โดยคำถามง่ายไม่ต้องค้น คำถามยากจึงค้นหลายรอบ \textbf{ส่วนที่คล้าย} คือมีการตัดสินใจว่าจะค้นหรือไม่ค้น \textbf{ส่วนที่ต่าง} คือใช้ตัวจำแนกที่ฝึกมาเฉพาะ ขณะที่โครงงานนี้ใช้แบบจำลองภาษาตัดสินใจโดยตรง \textbf{ข้อสังเกตที่นำมาใช้} คืองานนั้นเตือนว่าการค้นวนในคำถามที่ไม่ต้องวนทำให้ผลแย่ลง ซึ่งตรงกับผลเชิงลบที่พบในโครงงานนี้

Yan และคณะ \cite{yan2024crag} เสนอกลไกตรวจสอบคุณภาพของสิ่งที่ค้นได้ก่อนนำไปเขียนคำตอบ และแก้ไขเมื่อพบว่าค้นได้ไม่ดี \textbf{ส่วนที่คล้าย} คือมีขั้นประเมินคุณภาพผลการค้น \textbf{ส่วนที่ต่าง} คืองานนั้นใช้ตัวประเมินแยกที่ฝึกมาเฉพาะ ส่วนโครงงานนี้ใช้เกณฑ์สัมบูรณ์จากคะแนนความใกล้เชิงความหมายร่วมกับจำนวนสินค้าที่ผ่านเกณฑ์

\subsection{งานทบทวนวรรณกรรมที่ใช้กำหนดนิยาม}

Gao และคณะ \cite{gao2024modular} จัดระบบ RAG เป็นโมดูลที่ประกอบสลับกันได้ และจัดให้การค้นวนรอบอยู่ในกลุ่ม RAG ขณะที่ Singh และคณะ \cite{singh2025agentic} ทบทวนเฉพาะ Agentic RAG และระบุว่าเกณฑ์คือการมีชั้นที่แบบจำลองตัดสินใจเอง \textbf{โครงงานนี้ใช้สองงานนี้เป็นเกณฑ์จัดประเภท}ให้กับสถาปัตยกรรมทั้งสามแบบ ไม่ได้ตั้งนิยามขึ้นเอง

\subsection{ชุดข้อมูลมาตรฐานด้านการค้นหาสินค้า}

Reddy และคณะ \cite{reddy2022esci} เผยแพร่ชุดข้อมูลคำค้นสินค้าจริงของผู้ใช้ Amazon พร้อมเกณฑ์ความตรง 4 ระดับ ส่วน Chen และคณะ \cite{chen2022wands} เผยแพร่ชุดข้อมูลลักษณะเดียวกันจากเว็บขายเฟอร์นิเจอร์ \textbf{โครงงานนี้ยืมข้อความคำค้นจากทั้งสองชุด แต่ไม่ได้ยืมเฉลย} เนื่องจากคลังสินค้าเป็นคนละคลังกัน เฉลยของทั้งสองชุดจึงใช้กับคลังของโครงงานนี้ไม่ได้ \textbf{ข้อจำกัดร่วมของทั้งสองชุด} คือไม่มีคำค้นภาษาไทย

\subsection{ตารางเปรียบเทียบงานที่เกี่ยวข้อง}

การเปรียบเทียบคุณสมบัติของงานที่เกี่ยวข้องกับโครงงานนี้แสดงดังตารางที่ \ref{tab:related_compare}

\begin{table}[H]
\centering
\caption{เปรียบเทียบงานวิจัยที่เกี่ยวข้องกับโครงงานนี้}
\label{tab:related_compare}
\begin{tabular}{|l|c|c|c|c|}
\hline
\textbf{คุณสมบัติ} & \textbf{RAG} & \textbf{Adaptive} & \textbf{CRAG} & \textbf{โครงงานนี้} \\
 & \cite{lewis2020rag} & \cite{jeong2024adaptive} & \cite{yan2024crag} & \\
\hline
ค้นซ้ำได้หลายรอบ & -- & \cmark & \cmark & \cmark \\ \hline
แบบจำลองภาษาตัดสินใจเองขณะทำงาน & -- & -- & -- & \cmark \\ \hline
รองรับคำค้นภาษาไทย & -- & -- & -- & \cmark \\ \hline
โดเมนสินค้าเพื่อการซื้อขาย & -- & -- & -- & \cmark \\ \hline
วัดการปฏิเสธเมื่อคลังไม่มีของ & -- & -- & -- & \cmark \\ \hline
เปลี่ยนตัวแปรทีละตัวในการเปรียบเทียบ & -- & -- & -- & \cmark \\ \hline
รันบนเครื่องส่วนบุคคลทั้งหมด & -- & -- & -- & \cmark \\ \hline
เผยแพร่ชุดทดสอบพร้อมเฉลย & \cmark & \cmark & \cmark & \cmark \\
\hline
\end{tabular}
\end{table}

\subsection{ช่องว่างของงานวิจัยที่โครงงานนี้เติมเต็ม}

จากการทบทวนงานที่เกี่ยวข้อง พบช่องว่าง 4 ประการ

\begin{enumerate}[label=\arabic*)]
    \item \textbf{การเปรียบเทียบที่เปลี่ยนตัวแปรทีละตัว} งานที่เสนอสถาปัตยกรรมใหม่มักเทียบกับ RAG พื้นฐานโดยเปลี่ยนหลายอย่างพร้อมกัน จึงบอกไม่ได้ว่าผลที่ดีขึ้นมาจากส่วนใด โครงงานนี้ออกแบบให้สถาปัตยกรรมที่นำมาเทียบใช้โมดูลร่วมกันทั้งหมด ต่างกันเฉพาะชั้นที่ตัดสินใจ
    \item \textbf{การประเมินในบริบทสองภาษาที่ข้อมูลเป็นภาษาเดียว} งานที่มีอยู่ประเมินบนคำค้นและเอกสารภาษาเดียวกัน ขณะที่สถานการณ์จริงของผู้ใช้ไทยคือคำค้นเป็นภาษาไทยแต่ข้อมูลสินค้าเป็นภาษาอังกฤษ
    \item \textbf{การวัดการปฏิเสธที่ถูกต้อง} ตัวชี้วัดมาตรฐานวัดเฉพาะคุณภาพของสิ่งที่ตอบ ไม่ได้ให้เครดิตกับระบบที่ตอบว่าไม่พบเมื่อคลังไม่มีของจริง ซึ่งเป็นพฤติกรรมที่ระบบแนะนำสินค้าต้องทำให้ถูก
    \item \textbf{ข้อจำกัดของแบบจำลองขนาดที่รันได้บนเครื่องทั่วไป} งานที่มีอยู่ส่วนใหญ่ใช้แบบจำลองขนาดใหญ่ผ่านบริการภายนอก จึงยังไม่มีหลักฐานว่าความสามารถในการตัดสินใจของแบบจำลองขนาดที่รันบนเครื่องส่วนบุคคลได้นั้นเพียงพอหรือไม่
\end{enumerate}
